\section{The actor}
\label{sec:actor}

APT29 is tracked under several names. Microsoft calls the group Midnight Blizzard or NOBELIUM, Mandiant used UNC2452 for the SolarWinds activity before merging it into APT29~\cite{mandiant-unc2452}, and older reporting uses Cozy Bear or The Dukes~\cite{mitre-apt29,eset-ghost}. A 2022 joint advisory from agencies in the United States, Australia, Canada, New Zealand, and the United Kingdom states that the SVR has run cyber operations since at least 2008 and lists APT29 among the names used for those actors~\cite{aa22-110a}.

ESET reported that Operation Ghost began in 2013 and was still active in 2019, targeting foreign ministries in at least three European countries and an EU member state's embassy in Washington~\cite{eset-ghost}. In the Democratic National Committee breach, CrowdStrike found Cozy Bear on the network from summer 2015 and a separate actor, Fancy Bear, from April 2016, with no evidence that the two worked together~\cite{crowdstrike-dnc}. In July 2020, the UK National Cyber Security Centre and Canada's Communications Security Establishment reported that APT29 had targeted organisations involved in COVID-19 vaccine development in Canada, the United States, and the United Kingdom, with the likely aim of obtaining vaccine research~\cite{ncsc-vaccine}.

After SolarWinds the group changed access methods without changing its target set. It exploited JetBrains TeamCity in 2023~\cite{aa23-347a}, then ran phishing campaigns against German political parties, European diplomats, and more than 100 organisations in 2024 and 2025~\cite{mandiant-wineloader,ms-rdp,checkpoint-2025}. Check Point called the 2025 attribution likely, citing overlap with earlier activity and similar code.
